% ============================================================================
% kapitel/q.tex
% ============================================================================

\section{Q: Die Quelle der Organisation}

Q ist das zentrale Prinzip der IWT.

Es ist nicht das Bohm-Potential.

Es ist die Organisation der Information.

\section{Die Definition von Q}

Q ist definiert als das Prinzip, das die Information so organisiert, dass sie stabil ist.

\[
Q = \text{Organisation}(I)
\]

Es ist die Quelle aller Struktur.

\section{Die Herleitung von Q}

Q folgt aus zwei Prinzipien:

\begin{enumerate}
    \item \textbf{Informationserhaltung:} \( \sum |I|^2 = \text{const} \)
    \item \textbf{Minimale Spannung:} Die Information wird so organisiert, dass die Änderung zwischen benachbarten Knoten minimal ist.
\end{enumerate}

Die Lagrange-Multiplikator-Methode ergibt:

\[
\Delta^2 I = \lambda \cdot I
\]

Daraus folgt Q als das Potential, das diese Bedingung erfüllt.

\section{Q als Quelle}

Q ist die Quelle von:

\begin{itemize}
    \item \textbf{DBT:} Q ist das Bohm-Potential
    \item \textbf{WG:} \( F = -\nabla Q \) (Gravitation)
    \item \textbf{WED:} \( F = -\nabla Q \) (Elektrodynamik)
\end{itemize}

Alles kommt aus Q.

\section{Q als Selbstorganisation}

Q ist das Prinzip der Selbstorganisation.

Die Information organisiert sich selbst.

Q ist der Name dieser Organisation.

Q ist die Quelle.